\makeatletter
\def\input@path{{./}{../}{../../}{preamble/}{../preamble/}{../../preamble/}}
\makeatother
\input{preamble_exams.tex}

\geometry{left=0.48in,right=0.48in,top=0.42in,bottom=0.48in}
\linespread{1.03}
\setlength{\parskip}{2pt}
\setlist[enumerate,itemize]{topsep=1pt,itemsep=1pt,parsep=0pt,leftmargin=*}
\setlength{\columnsep}{0.25in}

\definecolor{algebrablue}{HTML}{24527A}
\newtcolorbox{keybox}{colback=blue!4,colframe=algebrablue,boxrule=.6pt,
  arc=1mm,left=2mm,right=2mm,top=1mm,bottom=1mm,before skip=3pt,after skip=3pt}
\newtcolorbox{examplebox}[1]{colback=black!2,colframe=black!45,boxrule=.5pt,
  arc=1mm,left=2mm,right=2mm,top=1mm,bottom=1mm,before skip=3pt,after skip=3pt,
  title=\textbf{#1},fonttitle=\small,coltitle=black}
\newcommand{\heading}[1]{\vspace{3pt}{\large\bfseries\color{algebrablue}#1}\par\vspace{1pt}}

\title{\vspace{-13mm}\textbf{Linear Equations in One Variable}\\[-1mm]
  \large Concise Learning and Practice Guide}
\author{}
\date{}
\pagestyle{empty}

\begin{document}
\maketitle
\vspace{-11mm}

\begin{multicols}{2}
\heading{1. What does solving mean?}
An \textbf{equation} states that two expressions are equal. A \textbf{variable}
(such as $x$) represents an unknown number. To \textbf{solve} an equation is to
find every value of the variable that makes the two sides equal.

\begin{keybox}
Think of an equation as a balance. To keep it balanced, perform the \emph{same
operation} on both sides. Equivalent steps preserve the solution.
\end{keybox}

\heading{2. Basic solving process}
Equations such as $ax+b=c$ and $ax+b=cx+d$ can be solved by:
\begin{enumerate}
  \item simplifying each side (distribute and combine like terms);
  \item collecting variable terms on one side;
  \item collecting constants on the other side;
  \item isolating the variable; and
  \item checking by substitution in the \emph{original} equation.
\end{enumerate}

\begin{examplebox}{Model: $3x+5=20$}
\vspace{-3mm}
\begin{align*}
3x+5-5&=20-5 &&\text{(subtract 5)}\\
3x&=15\\
\frac{3x}{3}&=\frac{15}{3} &&\text{(divide by 3)}\\
x&=5.
\end{align*}
\vspace{-4mm}
Check: $3(5)+5=20$, so the left side equals the right side. Thus $x=5$.
\end{examplebox}

\heading{3. Worked examples}
\begin{examplebox}{A. Form $ax+b=c$: $4x-7=17$}
\vspace{-3mm}
\begin{align*}
4x-7+7&=17+7\\4x&=24\\x&=6.
\end{align*}
\vspace{-4mm}
Check: $4(6)-7=17$.
\end{examplebox}

\begin{examplebox}{B. Variables on both sides: $5x+4=2x+19$}
\vspace{-3mm}
\begin{align*}
5x-2x+4&=2x-2x+19\\3x+4&=19\\
3x&=15\\x&=5.
\end{align*}
\vspace{-4mm}
Check: $5(5)+4=29=2(5)+19$.
\end{examplebox}

\begin{examplebox}{C. Parentheses: $3(x+2)=18$}
\vspace{-3mm}
\begin{align*}
3x+6&=18\\3x&=12\\x&=4.
\end{align*}
\vspace{-4mm}
Check: $3(4+2)=18$.
\end{examplebox}

\begin{examplebox}{D. Fraction: $\dfrac{x}{4}+3=8$}
\vspace{-3mm}
\begin{align*}
\frac{x}{4}&=5\\4\left(\frac{x}{4}\right)&=4(5)\\x&=20.
\end{align*}
\vspace{-4mm}
Check: $20/4+3=8$.
\end{examplebox}

\heading{4. Equations in context}
Follow the chain
\[
\text{context}\longrightarrow\text{equation}\longrightarrow
\text{solution}\longrightarrow\text{interpretation}.
\]

\textbf{Taxi.} A taxi charges a \$4 fixed fee and \$2 per kilometre. The fare is
\$18. Let $k$ be the kilometres travelled.
\[
4+2k=18\;\Longrightarrow\;2k=14\;\Longrightarrow\;k=7.
\]
The trip was \textbf{7 km}.

\textbf{Online order.} Four identical notebooks plus a \$3 delivery charge cost
\$27. Let $p$ be the price of one notebook.
\[
4p+3=27\;\Longrightarrow\;4p=24\;\Longrightarrow\;p=6.
\]
Each notebook costs \textbf{\$6}.

\textbf{Saving.} Maya already has \$25 and saves \$8 each week. When will she
have \$81? Let $w$ be the number of weeks.
\[
25+8w=81\;\Longrightarrow\;8w=56\;\Longrightarrow\;w=7.
\]
Maya needs \textbf{7 weeks}.

\heading{5. Common mistakes}
\begin{itemize}
  \item \textbf{Changing one side only:} from $x+3=9$, write
  $x+3-3=9-3$, not $x=9$.
  \item \textbf{Losing a negative sign:} $-2x=10$ gives $x=-5$, not $5$.
  \item \textbf{Distributing incorrectly:} $-3(x-2)=-3x+6$, not $-3x-6$.
  \item \textbf{Combining unlike terms:} $3x+4$ cannot become $7x$.
  \item \textbf{Stopping at a number in context:} state its meaning and units,
  for example ``7 kilometres,'' not just ``7.''
\end{itemize}
\end{multicols}

\newpage
\vspace*{-7mm}
\begin{multicols}{2}
\heading{6. Practice problems}
Solve each equation. Check your answer when possible.
\begin{enumerate}[label=\textbf{\arabic*.}]
  \item $x-9=6$
  \item $7x=42$
  \item $3x+4=25$
  \item $6x-5=31$
  \item $9-2x=19$
  \item $7x+3=4x+24$
  \item $8x-11=3x+9$
  \item $5(x-2)=30$
  \item $-2(3x+1)=16$
  \item $3(x+4)-2=2x+9$
  \item $\dfrac{x}{5}-2=4$
  \item $\dfrac{2x}{3}+1=9$
\end{enumerate}

For Problems 13--16, \textbf{define a variable, construct an equation, solve,
and interpret} the answer.
\begin{enumerate}[label=\textbf{\arabic*.},start=13]
  \item A delivery service charges a fixed fee of \$5 plus \$3 per kilometre.
  A delivery costs \$26. How far was it delivered?
  \item Six identical movie tickets and a \$4 booking fee cost \$58. What was
  the price of one ticket?
  \item Leo has \$18 and saves \$7 each week. How many weeks will it take him
  to have \$67?
  \item A school club paid \$96 for a room plus 8 identical meal packages. The
  total bill was \$216. What did each meal package cost?
\end{enumerate}

\heading{7. Challenge problems}
\begin{enumerate}[label=\textbf{C\arabic*.}]
  \item $4(2x-3)+5=3(x+6)$
  \item $\dfrac{x-2}{3}+\dfrac{x+1}{2}=7$
  \item A gym offers Plan A for \$18 per month plus a one-time \$24 joining fee.
  Plan B costs \$22 per month with no joining fee. After how many months do the
  plans have the same total cost? Define a variable, write an equation, solve,
  and interpret.
\end{enumerate}

\columnbreak
\heading{8. Answer key}
\textbf{Practice}
\begin{enumerate}[label=\textbf{\arabic*.},itemsep=2pt]
  \item $x=15$
  \item $x=6$
  \item $x=7$
  \item $x=6$
  \item $x=-5$
  \item $x=7$
  \item $x=4$
  \item $x=8$
  \item $x=-3$
  \item $x=-1$
  \item $x=30$
  \item $x=12$
  \item Let $k$ be the distance in kilometres. $5+3k=26$; $k=7$. The delivery
  distance was \textbf{7 km}.
  \item Let $p$ be one ticket's price. $6p+4=58$; $p=9$. Each ticket cost
  \textbf{\$9}.
  \item Let $w$ be the number of weeks. $18+7w=67$; $w=7$. Leo needs
  \textbf{7 weeks}.
  \item Let $m$ be one meal package's price. $96+8m=216$; $m=15$. Each package
  cost \textbf{\$15}.
\end{enumerate}

\textbf{Challenge}
\begin{enumerate}[label=\textbf{C\arabic*.},itemsep=2pt]
  \item $x=5$
  \item $x=\dfrac{43}{5}$
  \item Let $m$ be the number of months. $18m+24=22m$; $m=6$. The plans cost
  the same after \textbf{6 months} (\$132 each).
\end{enumerate}

\begin{keybox}
\textbf{Final check:} Does your value make the original left and right sides
equal? In a word problem, is the value reasonable, and have you included its
meaning and units?
\end{keybox}
\end{multicols}
\end{document}
